\section{Model}\label{sec.model}

\paragraph{Parties and corruption}

\begin{assumption}[Static corruption]\label{asm.static-corruption}
There are $n$ parties.
The adversary corrupts a fixed set of at most $t$ parties before the execution starts.
\end{assumption}

\begin{definition}[Honest and corrupted parties]\label{def.honest-corrupted}
The corrupted parties are the parties in the set of \Cref{asm.static-corruption}.
The other parties are honest.
\end{definition}

\begin{assumption}[Corruption bounds]\label{asm.corruption-bounds}
At least two parties are honest: $n - t \geq 2$.
The adversary may corrupt at least one party: $t \geq 1$.
\end{assumption}

\paragraph{Rounds}

\begin{assumption}[Rounds]\label{asm.rounds}
The execution proceeds in rounds $r = 1, 2, 3, \dots$.
Every party knows the current round.
\end{assumption}

\begin{assumption}[Genesis]\label{asm.genesis}
Before round $1$, that is, in round $0$, every party holds the genesis chain $\genesis$.
\end{assumption}

\paragraph{Chains}

\begin{definition}[Chain and length]\label{def.chain-length}
A chain is a sequence of blocks that starts with $\genesis$.
The length $\len{C}$ of a chain $C$ is the number of blocks of $C$ after $\genesis$.
So $\len{\genesis} = 0$.
\end{definition}

\begin{assumption}[Valid chains]\label{asm.valid-chains}
Only valid chains exist.
The adversary cannot forge blocks.
Every chain an honest party receives is valid.
\end{assumption}

\paragraph{Honest round procedure}

\begin{assumption}[Honest round procedure]\label{asm.honest-round}
In round $r$, an honest party $P$ executes the following steps, in this order.
\begin{enumerate}
    \item Receive every chain delivered to $P$ in round $r$.
    \item Adopt the longest chain among its current chain and the received chains.
          On a tie, keep the current chain.
    \item Make one mining attempt.
          On success, extend the current chain by one block.
    \item If the current chain changed during round $r$, broadcast it.
\end{enumerate}
\end{assumption}

\paragraph{Mining}

\begin{assumption}[Honest mining attempts]\label{asm.honest-mining}
Each honest party makes one mining attempt per round.
\end{assumption}

\begin{assumption}[Corrupted mining attempts]\label{asm.corrupted-mining}
Each corrupted party makes one mining attempt per round, on any chain the adversary chooses.
\end{assumption}

\begin{assumption}[Success patterns]\label{asm.success-patterns}
Whether a mining attempt succeeds is part of the execution.
Every pattern of successes is a possible execution.
\end{assumption}

\begin{definition}[Successful round]\label{def.successful-round}
Round $r$ is successful if at least one honest mining attempt succeeds in round $r$.
\end{definition}

\paragraph{Network}

\begin{assumption}[Honest delivery]\label{asm.honest-delivery}
A chain that an honest party broadcasts in round $r$ is delivered to every honest party in round $r + 1$.
\end{assumption}

\begin{assumption}[No dropping or delaying]\label{asm.no-drop}
The adversary cannot drop or delay the deliveries of \Cref{asm.honest-delivery}.
\end{assumption}

\begin{assumption}[Adversarial delivery]\label{asm.adversarial-delivery}
In any round, the adversary may deliver any chain it knows to any subset of the honest parties, and to none of the others.
\end{assumption}

\begin{assumption}[Rushing]\label{asm.rushing}
The adversary is rushing: in round $r$, it sees the honest broadcasts of round $r$ before it chooses its own deliveries for round $r + 1$.
\end{assumption}

\paragraph{Notation}

\begin{definition}[Chain of an honest party]\label{def.chain-at}
The chain of honest party $P$ at the end of round $r$ is $\chainat{P}{r}$, and $\chainat{P}{0} = \genesis$.
\end{definition}

\begin{definition}[Minimum and maximum length]\label{def.min-max-len}
The minimum of $\len{\chainat{P}{r}}$ over all honest $P$ is $\minlen{r}$, and the maximum is $\maxlen{r}$.
So $\minlen{0} = \maxlen{0} = 0$.
\end{definition}
